\chapter*{Glossary}
\addcontentsline{toc}{chapter}{Glossary}

The vocabulary the manual uses, in alphabetical order. Terms that
appear in code identifiers carry the typewriter font of the
identifier as well as the prose definition.

\begin{description}
  \item[Allocated size.] Bytes the filesystem has assigned to a
    file's content. Distinct from logical size (the user-visible
    byte length) and from exclusive size (bytes only this inode
    owns). Out of v1 scope; opens as R2-A.
  \item[APSB.] APFS volume superblock
    (\texttt{apfs\_superblock\_t}). Reached through the container
    OMAP at the scan XID; carries the volume OMAP, the FS-tree
    root, feature masks, role, UUID, and statistics.
  \item[Block zero.] The fixed locator block at the start of an
    APFS container. Carries enough metadata to find the
    checkpoint descriptor area but is not itself authoritative for
    the live container state.
  \item[Checkpoint descriptor area.] A contiguous run of blocks
    that hold the rotating ring of container superblocks and
    their associated checkpoint maps. The authoritative live NXSB
    lives here.
  \item[Checkpoint map.] A block carrying ephemeral-OID mappings
    for the selected checkpoint. Walked in order from
    \texttt{nx\_xp\_desc\_index}; the last block carries the
    \texttt{CHECKPOINT\_MAP\_LAST} flag.
  \item[Directory aggregate.] A per-directory total expressed
    under the unique-inode logical-sum policy: each inode
    contributes its logical size at most once per aggregate
    root, regardless of how many hard-linked paths inside that
    root name it.
  \item[\texttt{DIR\_REC}.] Directory record. The FS-tree leaf that
    represents a directory entry; key carries the parent inode's
    file ID and the name; value carries the child file ID and the
    POSIX entry type.
  \item[Dstream.] An APFS data stream. Logical and allocated size
    plus a default-extent reference; carried inline in the inode's
    xfields as \texttt{INO\_EXT\_TYPE\_DSTREAM} or in a separate
    \texttt{j\_dstream\_id\_val\_t} record.
  \item[Ephemeral object.] An APFS object that lives in memory
    between transactions and is written into the checkpoint data
    area at commit. The checkpoint map identifies it by
    \texttt{(oid, paddr)}.
  \item[Fail closed.] The project's discipline of refusing
    unknown, malformed, or out-of-allowlist input with a typed
    reason, rather than producing a partial or best-effort
    result.
  \item[Fallback path.] The supported-API backend that runs when
    raw mode rejects a source. POSIX walk plus \texttt{lstat}
    plus optional \texttt{getattrlistbulk}; emits the same row
    shape as raw mode.
  \item[Firmlink.] A macOS presentation mechanism that joins
    selected paths from a System volume to corresponding paths on
    a Data volume to produce the user-visible startup namespace.
    Out of v1 scope.
  \item[\texttt{FsRecordDump} / \texttt{FsRecordRow}.] The Rust
    crate's structured output for the FS-tree record-family dump
    and one decoded record respectively. Field-level equal to the
    independent Python decoder on the proof fixture.
  \item[FS tree.] The volume-level B-tree of variable-length
    records keyed by \texttt{(object\_id, type)}. Holds directory
    entries, inodes, xattrs, sibling-link records, dstream
    records, and extents.
  \item[\texttt{getattrlistbulk}.] macOS syscall that returns
    metadata for many directory entries per call. Used by the
    fallback walker for the bulk perf path.
  \item[Hard link.] A directory entry that names an inode already
    named by another directory entry. APFS represents the
    relationship through SIBLING\_LINK and SIBLING\_MAP records;
    the unique-inode aggregate policy ensures the inode's logical
    size is counted once per aggregate root.
  \item[\texttt{INODE}.] Inode record. The FS-tree leaf that
    carries file/dir/symlink identity, mode, parent ID, link
    count, internal flags, and an xfield tail.
  \item[Logical size.] User-visible byte length. The v1 size
    metric, equal to POSIX \texttt{st\_size} for regular files
    and to the symlink target byte length for symlinks. Distinct
    from allocated, exclusive, shared, and snapshot-retained
    sizes.
  \item[Max XID.] The XID argument to an OMAP lookup. Returns the
    entry with the highest stored XID less than or equal to the
    requested value. The whole point of pinning the scan XID.
  \item[\texttt{NamespaceEntry}.] One row of the namespace
    output: path, entry kind, file ID, logical size, and
    optional symlink target. Stable across raw and fallback
    backends.
  \item[NX / NXSB.] APFS container superblock
    (\texttt{nx\_superblock\_t}). The block-zero copy is the
    locator; the authoritative selected NXSB lives in the
    descriptor area at the highest valid checkpoint XID.
  \item[OID.] Object identifier. A virtual address into an APFS
    OMAP's name space. Not by itself a safe cache key
    (chapter~10).
  \item[OMAP.] Object map. The B-tree that resolves virtual OIDs
    to physical addresses under a chosen scan XID. Container and
    volume each have their own.
  \item[Paddr.] Physical block address. The on-disk location of an
    APFS object. Returned by an OMAP lookup for virtual objects;
    equal to \texttt{o\_oid} for physical objects.
  \item[Parsable.] Source property: checkpoint and root discovery
    succeed; raw traversal runs. Necessary but not sufficient
    for "supported."
  \item[Raw mode.] The native APFS-parsing backend. Used when the
    source is in the raw-mode allowlist (detached \texttt{.dmg},
    caller-supplied raw device, explicitly stable APFS source).
  \item[Readable.] Source property: raw bytes can be opened and
    sampled. Cheap. Not, on its own, evidence of correctness.
  \item[Scan XID.] The transaction identifier the parser pins for
    the duration of a run. All object resolution is performed
    against this XID. Picked as the highest checksum-valid NXSB
    XID in the descriptor area.
  \item[Sibling link / sibling map.] Hard-link records.
    \texttt{j\_sibling\_link\_val\_t} carries a hard-link's
    sibling-group identifier and the linked name;
    \texttt{j\_sibling\_map\_val\_t} closes the mapping from
    sibling ID to inode.
  \item[Sparse bytes.] \texttt{INO\_EXT\_TYPE\_SPARSE\_BYTES}: the
    count of bytes the sparse-allocation layer has not committed
    to disk. An allocation hint, not a logical size. Treat as
    diagnostic in chapter~8's precedence rule.
  \item[Supported.] Source property: the source is in the product
    allowlist for the requested mode. Validation has been
    performed; results were correct; the source class belongs in
    production.
  \item[Validatable.] Source property: parser output can be
    compared against a named oracle for the requested mode. A
    prerequisite of "supported" but not, by itself, sufficient.
  \item[Virtual object.] An APFS object whose
    \texttt{obj\_phys\_t} storage flag indicates the object
    is resolvable through an OMAP. \texttt{o\_oid} is independent
    of the block's paddr; the header is validated against
    \texttt{omap\_lookup.oid}.
  \item[XID.] Transaction identifier. Monotonic per container.
    Used to order checkpoints (chapter~3) and to disambiguate OMAP
    entries (chapter~4).
  \item[Xfield.] The variable-length tail of an FS-tree record
    that carries extended fields. One source-backed cursor rule
    (chapter~6); fail-closed on any malformed shape
    (chapter~7).
\end{description}
